\documentclass[conference]{IEEEtran}

\usepackage{amsmath,amssymb,amsthm}
\usepackage{booktabs,array,tabularx,microtype,needspace,flushend,url,placeins,float}
\usepackage[hidelinks]{hyperref}
\hypersetup{
pdftitle={TIDE: A Contract for Comparing Exact Execution Relationships Across Execution Records},
pdfauthor={Gabriel Dolbear},
pdfkeywords={execution records, mapped-scope conformance, exact execution relationships, provenance interoperability}
}

\newtheoremstyle{tideproposition}{0.35ex}{0.35ex}{\normalfont}{}{\bfseries}{.}{0.45em}{}
\theoremstyle{tideproposition}
\newtheorem{proposition}{Proposition}
\newcommand{\direct}[3]{#1\xrightarrow{#2}#3}
\newcommand{\reach}{\rightsquigarrow}
\newcommand{\reviewitem}[1]{\par\addvspace{0.22\baselineskip}\noindent\emph{#1.}\ }
\newcolumntype{Y}{>{\raggedright\arraybackslash}X}
\newcommand{\same}{\ensuremath{=}}
\newcommand{\diff}{\ensuremath{\ne}}

\title{TIDE: A Contract for Comparing Exact Execution Relationships Across Execution Records}

\author{\IEEEauthorblockN{Gabriel Dolbear}
\IEEEauthorblockA{Independent Researcher\\\href{mailto:gabrieldolbear@gmail.com}{gabrieldolbear@gmail.com}}}

\begin{document}
\maketitle

\setlength{\abovedisplayskip}{0.9ex plus 2pt minus 1pt}
\setlength{\belowdisplayskip}{0.9ex plus 2pt minus 1pt}
\setlength{\abovedisplayshortskip}{0pt plus 2pt}
\setlength{\belowdisplayshortskip}{0.9ex plus 2pt minus 1pt}

\begin{abstract}
Systems retain \emph{execution records} describing completed work and the data it consumed or produced. These records support questions about origin, ancestry, and downstream use, but records of the same work may preserve different exact relationships. This paper audits 35 questions from the Second and Third Provenance Challenges and ML Metadata (MLMD), identifying six recurring requirement types. The Producer--Consumer relationship appears in 25 questions across all three corpora and is selected because it recurs, is domain-independent, and can be represented without replacing host-specific semantics. Its questions normalise to four answer forms: the producer and consumers of a data item, and the inputs and outputs of a process occurrence. Reviewed approaches can represent these answers, but admissible records do not uniformly retain enough information to determine all four. The review did not identify a shared contract that both requires their determination and provides a common basis for comparison. TIDE provides such a contract through explicit Execution and State identities, separate input and output incidence, validity rules, fixed handoff and reachability judgments, and comparison under supplied identity maps. A reference comparator applies the contract to Dagster and OpenLineage projections of the same six-operation workload. Under the supplied maps, the projections conform and have complete agreement; controlled cases separate partial scope, mapped disagreement, absent identities, and semantic invalidity. These results begin after host-record interpretation and identity correspondence are supplied. They do not establish that a projection is true or complete or that the correspondence is correct.
\end{abstract}

\begin{IEEEkeywords}
execution records, mapped-scope conformance, exact execution relationships, provenance interoperability
\end{IEEEkeywords}

\section{Introduction}
Execution records are retained by workflow systems, lineage event models, metadata stores, and portable provenance packages for different purposes and at different levels of detail~\cite{openlineage-object,mlmd-guide,wfrun}. When multiple records concern the same work, operators may still need to determine whether the represented relationships agree~\cite{mediator,ppif}.

This comparison is not byte-for-byte equality or schema unification. Identifiers, terminology, granularity, and host-specific meaning may differ across systems~\cite{wfrun,mediator}. Comparison instead requires a common answer, correspondence between the occurrences being compared, and a rule for determining whether the answers agree. Differently structured records may preserve the same answer, while similar-looking records may not.

This paper examines the exact answers requested by a bounded question corpus and identifies the recurring information requirements needed to produce them~\cite{artifact}. It selects one domain-independent requirement and reviews whether existing approaches represent, require, and compare it. TIDE then formalises the four answer forms supported by that requirement and defines their comparison under supplied occurrence correspondence. It distinguishes disagreement within mapped scope from information outside that scope and from mapped identities absent from a projection. Its contribution is a contract for determining and comparing those answers, not a workflow system or provenance vocabulary.

\Needspace{13\baselineskip}
\section{Research Questions}

\textbf{RQ1.} What exact answers are requested, and which recurring information requirements are needed to produce them?

\textbf{RQ2.} Which recurring requirement is selected for formalisation, how is it represented in existing approaches, do admissible records require its exact answers, and what do existing comparison mechanisms compare?

\textbf{RQ3.} What contract preserves the selected requirement, and what conditions define validity under that contract?

\textbf{RQ4.} What can one comparator determine about agreement across projections?

\section{Question Audit and Recurring Requirement Decomposition}
\subsection{Bounded Corpus and Audit Method}
The bounded corpus contains nine questions analysed from the Second Provenance Challenge, all four core and fifteen optional Third Challenge questions, and seven MLMD questions---six from the guide and one from the tutorial~\cite{mlmd-guide,second-challenge,third-challenge-questions,mlmd-tutorial}. These sources state concrete questions expected of provenance or lineage tools; MLMD adds a modern execution-metadata setting. The corpus closes when every question in those declared sets has been reviewed. The retained audit records the selection rationale, protocol, question mappings, and aggregation~\cite{artifact}.

For each question, the audit records: (1) the requested answer; (2) the target, type, and filter conditions already supplied; and (3) the information that must be preserved to produce that answer. Table~\ref{tab:need} gives the six resulting requirement types and their counts. One author performed these classifications; they are not source terminology, and the retained question-level map makes each decision inspectable.

\begin{table}[H]
\caption{Requirement types identified by the 35-question audit. Counts overlap because one question may require several types.}
\label{tab:need}
\centering
\footnotesize
\setlength{\tabcolsep}{2.4pt}
\begin{tabularx}{\columnwidth}{@{}p{0.25\columnwidth}Yp{0.11\columnwidth}@{}}
\toprule
Requirement type & Required when & Count\\
\midrule
Identity & exact recorded data items or process occurrences must be distinguishable. & 32\\
Producer--Consumer relationship & which process produced or consumed a data item, or which data items a process consumed or produced, must be known & 25\\
Execution Status & success, failure, completion, or halt state affects the answer & 9\\
Ordering/Time & time, elapsed duration, or order affects the answer & 6\\
Domain Semantics & domain-specific records, values, joins, queries, effects, or operation meaning must be interpreted & 9\\
Annotations and Filtering & names, types, roles, parameters, metadata, or context select the answer & 23\\
\bottomrule
\end{tabularx}
\end{table}

\subsection{Selected Requirement and Boundary}
Each recurring type was checked against three conditions: it had to (1) recur across questions; (2) answer domain-independent questions; and (3) be representable without replacing domain-specific semantics. Producer--Consumer relationship meets all three, appearing in 25 questions across the Second Challenge, Third Challenge, and MLMD.

The other types remain necessary where recorded: Identity distinguishes participants; Annotations and Filtering select them; Execution Status and Ordering/Time qualify answers; and Domain Semantics supplies source-specific interpretation. They are not selected for formalisation here.

\subsection{Exact Answer Forms and Required Distinctions}
The selected Producer--Consumer questions normalise to four exact answer forms~\cite{artifact}:

\begin{table}[H]
\caption{Selected answer forms and the distinctions required to state them exactly.}
\label{tab:answer-distinctions}
\centering
\footnotesize
\setlength{\tabcolsep}{2.5pt}
\begin{tabularx}{\columnwidth}{@{}Y Y@{}}
\toprule
Exact answer & Required distinctions\\
\midrule
Which process occurrence produced a selected data item? & data-item occurrence identity; process-occurrence identity; production direction\\
Which process occurrences consumed a selected data item? & data-item occurrence identity; process-occurrence identity; consumption direction\\
Which exact data items did a process occurrence consume? & process-occurrence identity; data-item occurrence identity; consumption direction\\
Which exact data items did a process occurrence produce? & process-occurrence identity; data-item occurrence identity; production direction\\
\bottomrule
\end{tabularx}
\end{table}

Across the four forms, three distinctions recur: process-occurrence identity, data-item occurrence identity, and production versus consumption direction. Section~V represents these distinctions; Section~VI-B erases each one.

Questions requiring the other types in Table~\ref{tab:need} remain outside the selected contract. Section~IV reviews the selected requirement in existing approaches.

\section{Existing Formalisations and the Remaining Gap}

\subsection{Review Questions}
Table~\ref{tab:prior} checks representability, admissible-record sufficiency, the comparison target, and what interpretation or correspondence must already be supplied. The review is scoped to the selected requirement and does not rank general provenance expressiveness.

\begin{table*}[!t]
\caption{Scoped review against the selected four-answer Producer--Consumer requirement. ``Required'' uses each work's own acceptance criterion; details follow in the text.}
\label{tab:prior}
\centering
\footnotesize
\renewcommand{\arraystretch}{1.12}
\setlength{\tabcolsep}{2.5pt}
\begin{tabularx}{\textwidth}{@{}
>{\raggedright\arraybackslash}p{0.09\textwidth}
>{\raggedright\arraybackslash}p{0.16\textwidth}
>{\raggedright\arraybackslash}p{0.19\textwidth}
>{\raggedright\arraybackslash}p{0.23\textwidth}
Y@{}}
\toprule
Work & Four answers retainable? & Required in every admissible record? & Comparison or cross-record check & Must be supplied\\
\midrule
OPM\newline\cite{opm} & \textbf{Yes}: typed Process--Artifact incidence. & \textbf{No}: \emph{was triggered by} may hide the Artifact identity. & OPM equality, profiles, persistent names, and refinement. & Profile conventions or cross-graph names.\\
PROV\newline\cite{prov-dm,prov-constraints} & \textbf{Yes}: typed Activity--Entity incidence. & \textbf{No}: \texttt{wasInformedBy} may leave only an existential Entity. & Normalisation and isomorphism of valid normal forms. & Resolved URI co-reference.\\
Mediator\newline\cite{mediator} & \textbf{Conditional}: wrappers must expose the relationships. & \textbf{No}: no source-independent completeness rule. & Common schema with local and transitive queries. & Wrapper semantics and identifiers.\\
CPF\newline\cite{cpf} & \textbf{Conditional}: the exchanged OPM graph must include them. & \textbf{No}: a dictionary reference may replace complete history. & Source IDs, naming dictionaries, and proposed Schematron checks. & Dictionary mappings.\\
ProvONE/\newline P-PIF\newline\cite{ppif} & \textbf{Conditional}: the retrospective trace must retain them. & \textbf{Not established} by the reviewed translation proof. & Common ProvONE/SPARQL representation and Prov2ONE translation. & Vocabulary rules and adapters.\\
Workflow\newline Run RO-Crate\newline\cite{wfrun,wfrun-profile} & \textbf{Conditional}: \texttt{object}/\texttt{result} and step detail must be present. & \textbf{No}: \texttt{object} is MAY; \texttt{result} is SHOULD. & Shared packaging and profile conformance. & Native identifiers and record interpretation.\\
ProvAbs\newline\cite{provabs} & \textbf{Conditional}: abstraction may erase exact nodes and edges. & \textbf{No}: permitted abstraction can remove exact identities. & Grouping transformation preserving stated validity conditions. & Grouping set and policy.\\
\bottomrule
\end{tabularx}
\end{table*}

\subsection{Can the Four Answers Be Represented?}
Yes, when the relevant detailed relationships are retained. OPM's \emph{used} and \emph{was generated by} edges and PROV's \texttt{used} and \texttt{wasGeneratedBy} relations directly record identified process--data incidence~\cite{opm,prov-dm}. The mediator's global schema and Workflow Run RO-Crate can likewise expose execution inputs and outputs when their source mappings or declared records provide them~\cite{mediator,wfrun}. The selected gap is therefore not whether the reviewed approaches can encode the four answers.

\subsection{Do Admissible Records Retain Enough?}
Not uniformly. OPM permits \emph{was triggered by}, where completion may establish a joining Artifact without retaining its identifier~\cite{opm}. PROV permits \texttt{wasInformedBy}; normalization can supply an existential Entity generated by one Activity and used by another without recovering the recorded occurrence identity required by the selected answers~\cite{prov-dm,prov-constraints}. CPF may export one Entity with a dictionary reference rather than its complete history~\cite{cpf}; Process Run Crate makes action \texttt{object} optional and \texttt{result} recommended rather than mandatory~\cite{wfrun-profile}; and ProvAbs can replace exact nodes during abstraction while preserving its stated validity conditions~\cite{provabs}. Thus an approach may be able to encode the four answers without requiring every admissible record to retain enough information to determine them.

\subsection{What Do the Comparison Mechanisms Check?}
OPM defines structural equality, persistent names, profiles, and refinement over inferred dependencies~\cite{opm}. PROV-CONSTRAINTS defines validity and equivalence through normalization and isomorphism of valid normal forms~\cite{prov-constraints}. CPF supplies source identifiers, naming dictionaries, and proposed Schematron checks~\cite{cpf}; mediation supplies a common query interface over source wrappers~\cite{mediator}; P-PIF supplies a common ProvONE representation and proves its prospective Prov2ONE translation~\cite{ppif}; and Workflow Run RO-Crate supplies shared packaging and run profiles for heterogeneous records~\cite{wfrun}. The review did not identify a shared contract that both requires complete determination of all four selected answer forms and provides a common basis for their comparison.

\Needspace{9\baselineskip}
\subsection{What Must Be Supplied Before Comparison?}
The reviewed approaches use application-level co-reference, identifiers, naming dictionaries, or wrapper mappings to supply identity alignment. PROV equivalence assumes that applications have already resolved URI co-reference~\cite{prov-constraints}. CPF uses source identifiers and naming dictionaries, while mediator wrappers determine how native fields, data items, and identifiers populate the common schema~\cite{cpf,mediator}. Automatic correspondence discovery is outside the scope of this paper.

\subsection{What Remains?}
The remaining obligation is a shared contract that determines the four selected answers and compares them under supplied identity correspondence. Section~V defines TIDE as one such contract.

\section{TIDE Exact-Relationship Contract}

\subsection{Candidate Structure and Identities}
A candidate TIDE structure is
\begin{equation}
K=(E,S,I,O),\qquad
E\subseteq\mathcal{E},\quad S\subseteq\mathcal{S},\quad
I,O\subseteq E\times S,
\label{eq:tide-structure}
\end{equation}
where $\mathcal{E}$ contains possible Execution identities and $\mathcal{S}$ contains possible State identities.

\emph{Execution.} An Execution identity denotes one occurrence of work, not a reusable job or operation type.

\emph{State.} A State identity denotes one represented data or condition occurrence. Two States may have equal values and still be different occurrences.

\emph{Input incidence.} $(e,s)\in I$ means Execution $e$ consumed State $s$.

\emph{Output incidence.} $(e,s)\in O$ means Execution $e$ produced State $s$.

The host supplies the represented occurrence identities $E$ and $S$; TIDE does not derive them from names, values, identifier syntax, or incidence. Explicit identity carriers distinguish a represented occurrence with an empty answer set from an occurrence absent from the structure. The relations $(I,O)$ are the relationship-bearing kernel: once $E$ and $S$ are fixed, all selected relationship answers are determined by input and output incidence.

\begin{table}[H]
\caption{Required distinctions and their TIDE representation.}
\label{tab:tide-representation}
\centering
\footnotesize
\setlength{\tabcolsep}{2.5pt}
\begin{tabularx}{\columnwidth}{@{}Y Y@{}}
\toprule
Required distinction & TIDE representation\\
\midrule
Process-occurrence identity & Execution identities in $E$\\
Data-item occurrence identity & State identities in $S$\\
Production versus consumption direction & Separate input incidence $I$ and output incidence $O$\\
\bottomrule
\end{tabularx}
\end{table}

\subsection{Fixed Relationship Judgments}
The notation $K\vdash J$ means that judgment $J$ is derived from the incidence in $K$.

\emph{Direct handoff.}
\begin{equation}
\frac{(e_p,s)\in O\qquad(e_c,s)\in I}
     {K\vdash\direct{e_p}{s}{e_c}}.
\end{equation}
A direct handoff exists when an Execution produced a State and an Execution consumed that same State.

\emph{Directed reachability.}
\begin{equation}
\frac{n\geq1,\ e_0=e_a,\ e_n=e_b,\ \forall i<n:\ K\vdash\direct{e_i}{s_i}{e_{i+1}}}
     {K\vdash e_a\reach e_b}.
\end{equation}
Directed reachability exists when one or more direct handoffs form a chain from $e_a$ to $e_b$. For fixed identity carriers, both judgments are determined only by $I$ and $O$. Neither asserts timing, physical transfer, value-level causal necessity, or capture completeness.

\subsection{Validity}

\subsubsection{Structural Validity}
A structurally valid TIDE representation has the form in Eq.~\ref{eq:tide-structure}: represented Execution and State identity sets with input and output incidence contained in their Cartesian product.
\subsubsection{Semantic Validity}
Semantic validity is separate from determination, comparison, and structural validity. It constrains the meaning carried by a structurally valid representation, not whether its answers can be computed. Fixed $E$, $S$, $I$, and $O$ still determine the selected answers, and a violation of either semantic rule remains comparable. Unique creation requires every represented State to have one represented origin, while acyclic ancestry requires reachability not to return to its starting Execution. Together, these rules give TIDE its intended meanings of single-origin State occurrences and one-way execution ancestry.

\emph{Unique creation.} Every represented State has exactly one represented producer:
\begin{equation}
\forall s\in S,\;\exists!e\in E:(e,s)\in O.
\label{eq:unique-creation}
\end{equation}
A producerless State still has a determinate producer answer: the empty set. Unique creation is therefore not required for determination, but it is required for semantic validity. Missing production incidence cannot distinguish a State intentionally introduced at the boundary from one whose producer is unknown or omitted. It cannot preserve that distinction either: if the same State is later represented as the output of an Execution, the structure retains no fact showing that it entered at the boundary. TIDE therefore does not treat a missing producer as a boundary marker. This rule does not claim that the native record is complete; it only rejects a supplied projection that leaves a represented State without a represented origin. A host may represent an external or unknown producer with a boundary Execution that produces the State.

Multiple producers for a single State also leave the producer answer determinate, but the representation no longer gives that State occurrence one represented origin. Exactly one producer therefore prevents multiple production attributions to the same State.

\emph{Acyclic ancestry.} No Execution reaches itself:
\begin{equation}
\forall e\in E,\qquad K\nvdash e\reach e.
\label{eq:acyclic-ancestry}
\end{equation}
Reachability remains determined in a cyclic structure, so acyclicity is likewise not required for determination. A cycle permits a directed handoff chain to return to its starting Execution, so reachability no longer represents one-way execution ancestry. Semantic validity therefore excludes cycles.

\subsection{Supplied Interpretation and Mapped-Scope Conformance}
A host record becomes a TIDE structure only after the host selects which native facts mean consumption and production and which occurrence identities represent the Executions and States. TIDE starts from the resulting $K=(E,S,I,O)$; it does not choose that interpretation.

Let $K_1=(E_1,S_1,I_1,O_1)$ and $K_2=(E_2,S_2,I_2,O_2)$. Two supplied bijections between selected identity subsets declare which Execution and State identities correspond:
\begin{equation}
\begin{aligned}
f_E&:D_E\overset{\sim}{\longrightarrow}R_E,
&D_E,R_E&\subseteq\mathcal{E},\\
f_S&:D_S\overset{\sim}{\longrightarrow}R_S,
&D_S,R_S&\subseteq\mathcal{S}.
\end{aligned}
\label{eq:supplied-identity-maps}
\end{equation}
The maps may cover only part of either structure, and a supplied map identity need not be present in the corresponding structure.

\emph{Mapped and present.} An identity is supplied by the map and present in the structure:
\begin{align}
E_1^m&=E_1\cap D_E, & E_2^m&=E_2\cap R_E,\\
S_1^m&=S_1\cap D_S, & S_2^m&=S_2\cap R_S.
\label{eq:mapped-present}
\end{align}

\emph{Outside.} An identity is present in the structure but has no supplied mapping:
\begin{align}
E_1^{\mathrm{out}}&=E_1\setminus D_E,
& E_2^{\mathrm{out}}&=E_2\setminus R_E,\\
S_1^{\mathrm{out}}&=S_1\setminus D_S,
& S_2^{\mathrm{out}}&=S_2\setminus R_S.
\label{eq:outside-identities}
\end{align}

\emph{Absent.} An identity is supplied by the map but is not present in the structure:
\begin{align}
E_1^{\mathrm{abs}}&=D_E\setminus E_1,
& E_2^{\mathrm{abs}}&=R_E\setminus E_2,\\
S_1^{\mathrm{abs}}&=D_S\setminus S_1,
& S_2^{\mathrm{abs}}&=R_S\setminus S_2.
\label{eq:absent-identities}
\end{align}

Mapped incidence contains incidences whose Execution and State identities
are mapped and present:
\begin{align}
I_i^m&=I_i\cap(E_i^m\times S_i^m),\\
O_i^m&=O_i\cap(E_i^m\times S_i^m),\qquad i\in\{1,2\}.
\label{eq:mapped-incidence}
\end{align}
The mapped structure is $K_i^m=(E_i^m,S_i^m,I_i^m,O_i^m)$.

Outside incidence contains the remaining incidences:
\begin{align}
I_i^{\mathrm{out}}&=I_i\setminus I_i^m,\\
O_i^{\mathrm{out}}&=O_i\setminus O_i^m.
\label{eq:outside-incidence}
\end{align}

The pair map applies the supplied Execution and State maps to each mapped incidence:
\begin{equation}
F:D_E\times D_S\overset{\sim}{\longrightarrow}R_E\times R_S,
\qquad
F(e,s)=(f_E(e),f_S(s)).
\end{equation}
For a set of mapped incidences, $F$ replaces each Execution and State identity with its supplied corresponding identity.

\emph{Preservation.}
Preservation reads forward from $K_1$ to $K_2$: every mapped incidence present in $K_1$ must occur at its mapped location in $K_2$. Preservation failures are
\begin{align}
\Delta_I^-&=F(I_1^m)\setminus I_2^m,\\
\Delta_O^-&=F(O_1^m)\setminus O_2^m.
\label{eq:preservation-deltas}
\end{align}

\emph{Reflection.}
Reflection reads from $K_2$ back to $K_1$: every mapped-scope incidence present in $K_2$ must be reflected by a corresponding incidence in $K_1$. Reflection failures are
\begin{align}
\Delta_I^+&=I_2^m\setminus F(I_1^m),\\
\Delta_O^+&=O_2^m\setminus F(O_1^m).
\label{eq:reflection-deltas}
\end{align}

The structures conform inside the supplied map exactly when preservation and reflection both hold:
\begin{equation}
F(I_1^m)=I_2^m,
\qquad
F(O_1^m)=O_2^m.
\label{eq:scoped-conformance}
\end{equation}
Equivalently, conformance holds exactly when all four difference sets are empty. Outside incidence and absent mapped identities are reported separately and are not used to determine conformance.

\emph{Complete agreement across projections.}
Conformance compares only mapped incidence. The projections agree completely when conformance holds and the supplied maps cover exactly the represented Execution and State identities on both sides:
\begin{equation}
E_1=D_E,\quad E_2=R_E,\quad
S_1=D_S,\quad S_2=R_S.
\label{eq:projection-coverage}
\end{equation}

\section{Formal Guarantees}
The propositions establish determination, loss under erasure of the three Section~III distinctions, and mapped-scope answer agreement.

\subsection{Determination}
\begin{proposition}[Determination]\label{prop:determination}
For fixed carrier sets $E$ and $S$ in any structurally valid $K=(E,S,I,O)$, the relations $I$ and $O$ determine the four selected answer sets for every represented identity and the complete direct-handoff and directed-reachability answer sets.
\end{proposition}
\noindent\emph{Proof.}
For any represented State $s$ and Execution $e$, the four selected answers are exactly
\begin{align*}
\text{producer of }s &: \{e'\in E\mid(e',s)\in O\},\\
\text{consumers of }s &: \{e'\in E\mid(e',s)\in I\},\\
\text{inputs of }e &: \{s'\in S\mid(e,s')\in I\},\\
\text{outputs of }e &: \{s'\in S\mid(e,s')\in O\}.
\end{align*}
A direct handoff is fixed by the same incidence:
\[
K\vdash\direct{e_p}{s}{e_c}
\Longleftrightarrow
(e_p,s)\in O\land(e_c,s)\in I.
\]
Directed reachability is a finite non-empty chain of those handoffs. With the represented identities $E$ and $S$ held fixed, fixed $I$ and $O$ therefore fix every selected answer and both derived judgment sets, including empty answer sets. \hfill$\square$

\subsection{The Three Required Distinctions}
Table~\ref{tab:tide-representation} maps the three required distinctions to their TIDE representations. Each witness begins with two representations that have different answers. It then erases one distinction, after which the representations become the same.

\reviewitem{Execution occurrence identity}\par\nobreak
Let $E=\{e_1,e_2,e_3,e_4\}$ and $S=\{s_1,s_2\}$, and let
\[
O=\{(e_1,s_1),(e_2,s_2)\}.
\]
In $K_1$, let
\[
I_1=\{(e_3,s_1),(e_4,s_2)\};
\]
in $K_2$, let
\[
I_2=\{(e_4,s_1),(e_3,s_2)\}.
\]
Then
\[
K_1\vdash\direct{e_1}{s_1}{e_3},
\qquad
K_2\vdash\direct{e_1}{s_1}{e_4}.
\]
Erase Execution identity by replacing $e_1,e_2,e_3,e_4$ with $e$, so $E'=\{e\}$. Both representations then reduce to
\[
\begin{aligned}
O_1'&=O_2'=\{(e,s_1),(e,s_2)\},\\
I_1'&=I_2'=\{(e,s_1),(e,s_2)\}.
\end{aligned}
\]
Both direct-handoff answer sets therefore reduce to
\[
\begin{aligned}
K_1'&\vdash\direct{e}{s_1}{e},
&\qquad K_1'&\vdash\direct{e}{s_2}{e},\\
K_2'&\vdash\direct{e}{s_1}{e},
&\qquad K_2'&\vdash\direct{e}{s_2}{e}.
\end{aligned}
\]

\reviewitem{State occurrence identity}\par\nobreak
Let $E=\{e_1,e_2,e_3\}$ and $S=\{s_1,s_2\}$, and let
\[
O=\{(e_1,s_1),(e_2,s_2)\}.
\]
In $K_1$, let
\[
I_1=\{(e_3,s_1)\};
\]
in $K_2$, let
\[
I_2=\{(e_3,s_2)\}.
\]
Then
\[
K_1\vdash\direct{e_1}{s_1}{e_3},
\qquad
K_2\vdash\direct{e_2}{s_2}{e_3}.
\]
Erase State identity by replacing $s_1,s_2$ with $s$, so $S'=\{s\}$.
Both representations then reduce to
\[
O_1'=O_2'=\{(e_1,s),(e_2,s)\},
\qquad
I_1'=I_2'=\{(e_3,s)\}.
\]
Both direct-handoff answer sets therefore reduce to
\[
\begin{aligned}
K_1'&\vdash\direct{e_1}{s}{e_3},
&\qquad K_1'&\vdash\direct{e_2}{s}{e_3},\\
K_2'&\vdash\direct{e_1}{s}{e_3},
&\qquad K_2'&\vdash\direct{e_2}{s}{e_3}.
\end{aligned}
\]

\reviewitem{Production versus consumption direction}\par\nobreak
Let $E=\{e_1,e_2\}$ and $S=\{s_1\}$. In $K_1$, let
\[
O_1=\{(e_1,s_1)\},
\qquad
I_1=\{(e_2,s_1)\};
\]
in $K_2$, let
\[
O_2=\{(e_2,s_1)\},
\qquad
I_2=\{(e_1,s_1)\}.
\]
Then
\[
K_1\vdash\direct{e_1}{s_1}{e_2},
\qquad
K_2\vdash\direct{e_2}{s_1}{e_1}.
\]
Erase production versus consumption direction by replacing the ordered pair $(I,O)$ with the undirected union $I\cup O$. Both representations then reduce to
\[
I_1\cup O_1
=
I_2\cup O_2
=
\{(e_1,s_1),(e_2,s_1)\}.
\]
This common result does not determine whether the original answer was
\[
\direct{e_1}{s_1}{e_2}
\qquad\text{or}\qquad
\direct{e_2}{s_1}{e_1}.
\]

\begin{proposition}[Specified erasures do not preserve answers]
\label{prop:erasures}
None of the three displayed erasures preserves the complete direct-handoff
and directed-reachability answer sets for all structurally valid TIDE
representations.
\end{proposition}
\noindent\emph{Proof.}
Each pair has different direct-handoff answers before its named erasure but becomes the same after that erasure. In each pair, the differing handoff also gives a differing one-step reachability answer. The erased representation therefore does not determine which original answers applied. \hfill$\square$

\subsection{Mapped-Scope Answer Agreement}
Mapped conformance equates mapped input and output incidence under the supplied pair map $F$. The following proposition shows that handoff and reachability answers then agree in mapped scope.

To state judgments only over represented identities, let
\begin{align*}
E_1^{\leftrightarrow}&=\{e\in E_1^m\mid f_E(e)\in E_2^m\},\\
S_1^{\leftrightarrow}&=\{s\in S_1^m\mid f_S(s)\in S_2^m\}
\end{align*}
denote the source-side mapped identities whose corresponding target identities are also present.

\begin{proposition}[Mapped conformance preserves and reflects scoped answers]\label{prop:conformance}
If $K_1$ and $K_2$ conform under the supplied maps, then for every $e_p,e_c,e_a,e_b\in E_1^{\leftrightarrow}$ and $s\in S_1^{\leftrightarrow}$,
\begin{align}
K_1^m\vdash\direct{e_p}{s}{e_c}
&\Longleftrightarrow
K_2^m\vdash\direct{f_E(e_p)}{f_S(s)}{f_E(e_c)},\\
K_1^m\vdash e_a\reach e_b
&\Longleftrightarrow
K_2^m\vdash f_E(e_a)\reach f_E(e_b).
\end{align}
\end{proposition}
\noindent\emph{Proof.}
Conformance gives $F(I_1^m)=I_2^m$ and $F(O_1^m)=O_2^m$. A direct handoff consists of an output incidence and an input incidence sharing the same State. Therefore a mapped handoff exists on one side exactly when the corresponding mapped handoff exists on the other. A reachability answer is a finite non-empty chain of handoffs; mapping each handoff gives the corresponding chain, and the inverse bijections give the reverse direction. \hfill$\square$

The proposition applies only to jointly present mapped identities; outside incidence and absent identities are not compared, and semantic validity is not required.

\section{Evaluation: What One Comparator Determines}

\begin{table*}[!t]
\caption{Evaluation results under the definitions in Section~V.}
\label{tab:evaluation}
\centering
\footnotesize
\renewcommand{\arraystretch}{1.04}
\setlength{\tabcolsep}{1.7pt}
\begin{tabularx}{\textwidth}{@{}
Y
>{\centering\arraybackslash}p{0.10\textwidth}
>{\centering\arraybackslash}p{0.065\textwidth}
>{\centering\arraybackslash}p{0.065\textwidth}
>{\centering\arraybackslash}p{0.13\textwidth}
>{\centering\arraybackslash}p{0.13\textwidth}
>{\centering\arraybackslash}p{0.07\textwidth}
>{\centering\arraybackslash}p{0.10\textwidth}
@{}}
\toprule
Comparison &
Semantic validity $K_1/K_2$ &
$\Delta I$ $-/+$ &
$\Delta O$ $-/+$ &
Outside $I/O$ $K_1;K_2$ &
Absent $E/S$ $K_1;K_2$ &
Conforms &
Complete agreement \\
\midrule

Dagster/OpenLineage retained projections
& Y/Y & 0/0 & 0/0 & 0/0; 0/0 & 0/0; 0/0 & Y & Y\\

Controlled structures, full identity coverage
& Y/Y & 0/0 & 0/0 & 0/0; 0/0 & 0/0; 0/0 & Y & Y\\

Controlled structures, partial identity coverage
& Y/Y & 0/0 & 0/0 & 3/3; 3/3 & 0/0; 0/0 & Y & N\\

Omitted mapped input
& Y/Y & 1/0 & 0/0 & 0/0; 0/0 & 0/0; 0/0 & N & N\\

Added mapped-scope input
& Y/Y & 0/1 & 0/0 & 0/0; 0/0 & 0/0; 0/0 & N & N\\

Omitted mapped output
& Y/N & 0/0 & 1/0 & 0/0; 0/0 & 0/0; 0/0 & N & N\\

Mapped output, target identity absent
& Y/Y & 0/0 & 1/0 & 0/0; 0/1 & 0/0; 0/1 & N & N\\

Added mapped-scope output
& Y/Y & 0/0 & 0/1 & 0/0; 0/0 & 0/1; 0/0 & N & N\\

Invalid outside branch
& Y/N & 0/0 & 0/0 & 0/0; 1/0 & 0/0; 0/0 & Y & N\\

\bottomrule
\end{tabularx}
\end{table*}

\subsection{Evaluation Questions}
The evaluation tests whether (1) each retained TIDE projection determines the four selected answer forms; (2) the independently produced Dagster and OpenLineage projections agree under supplied occurrence correspondence; and (3) controlled changes are distinguished as preservation failure, reflection failure, outside incidence, absent mapped identity, semantic invalidity, or incomplete identity coverage.

\subsection{Evaluation Material and Reported Results}
Table~\ref{tab:evaluation} reports the retained and controlled comparison results.
One six-operation workload is executed twice: once by Dagster and once by a Python runner emitting OpenLineage events. Each native record is projected to $K=(E,S,I,O)$ and the identity maps are supplied externally. Controlled rows vary only the supplied map, represented identity carriers, or one incidence; they perform no further execution.

The reference comparator receives only $K_1$, $K_2$, and the two supplied identity maps. It checks map one-to-one form, derives the fixed judgments and semantic-validity result for each structure, and applies the mapped-comparison definitions from Section~V-D.

\subsection{Four Selected Answers in the Dagster/OpenLineage Records}
Both retained projections independently determine the four Section~III answer forms. The labels below denote corresponding mapped occurrences:
\begin{align*}
\operatorname{producer}(\texttt{merged})
   &= \{\texttt{merge}\},\\
\operatorname{consumers}(\texttt{merged})
   &= \{\texttt{analyse},\texttt{archive}\},\\
\operatorname{inputs}(\texttt{merge})
   &= \{\texttt{ready\_left},\texttt{ready\_right}\},\\
\operatorname{outputs}(\texttt{merge})
   &= \{\texttt{merged}\}.
\end{align*}

\subsection{Dagster/OpenLineage Cross-Record Comparison}
Each projection contains six Executions, seven States, six input incidences, and seven output incidences; each derives six State-labelled handoffs and thirteen non-empty reachability answers, and both are semantically valid. All mapped differences, outside incidence, and absent identities are empty, so the supplied correspondence yields conformance and complete agreement across projections.

\subsection{Controlled Comparison Cases}
The controlled rows begin from matching structures under different identities. Partial coverage changes only scope: conformance holds while complete agreement does not. Removing or adding mapped incidence produces the corresponding preservation or reflection failure. Absent identities are reported separately, while the producerless outside branch changes semantic validity without changing mapped conformance.

\subsection{Reproduction and Evaluation Boundary}
The retained artifact contains the comparator, controlled cases, native
events, projections, maps, expected counts, and regression tests. The evidence
does not establish that a supplied projection is true or complete or that
supplied occurrence correspondence is semantically correct.

\FloatBarrier
\section{Discussion and Limitations}
\subsection{Answers to the Research Questions}

\reviewitem{RQ1---Requested answers and recurring requirements}
Table~\ref{tab:need} reports six recurring requirement types. The selected
Producer--Consumer requirement reduces to four exact answer forms, requiring
process-occurrence identity, data-item occurrence identity, and production
versus consumption direction.

\reviewitem{RQ2---Selected requirement and existing approaches}
The reviewed approaches can represent the four answers when the required information is present, but admissible records do not uniformly require it. Their comparison mechanisms target structures, equivalence, translations, profiles, or common query representations. The review did not identify a shared contract that both requires all four answers and provides a common basis for their comparison; identity alignment remains supplied.

\reviewitem{RQ3---Contract and validity}
TIDE represents the required distinctions through Execution and State identities
and separate input and output incidence; once $E$ and $S$ are fixed, $I$ and $O$ determine the four
selected answers. Structural validity requires that representation form, while
semantic validity requires unique production attribution and acyclic execution ancestry.

\Needspace{7\baselineskip}
\reviewitem{RQ4---Agreement across projections}
Given supplied Execution and State maps, the comparator reports semantic validity, mapped scope, outside incidence, absent identities, preservation/reflection failures, conformance, and complete agreement. Proposition~\ref{prop:conformance} gives the corresponding handoff and reachability guarantee within conforming mapped scope.

\subsection{Explicit Boundaries of the Result}

\reviewitem{Supplied before TIDE}
The host selects and projects native records, assigns identities,
and supplies cross-record identity correspondence; TIDE does not discover or verify those choices.

\reviewitem{Compared and not compared}
Only mapped-and-present incidence is compared. Present but unmapped incidence
is outside; mapped but unrepresented identities are absent. Neither is inferred
or repaired.

\reviewitem{Relationship meaning}
Handoff records production and consumption of the same State; reachability is a non-empty chain of handoffs. Neither asserts timing, physical transfer, causal necessity, or capture completeness.

\reviewitem{Representation limits}
The supplied identity maps do not represent many-to-one correspondence.

\reviewitem{Evidence limits}
The retained pair evaluates one six-operation workload executed through Dagster and OpenLineage; controlled cases vary scope, incidence, identity presence, and semantic validity. They do not evaluate other domains or larger histories.

\reviewitem{Requirements outside the selected contract}
Filtering, status, ordering/time, and domain-semantic answers require their
additional representations or semantics.

\Needspace{10\baselineskip}
\section{Conclusion}
The 35-question audit selects Producer--Consumer relationship, which normalises to four answer forms requiring process-occurrence identity, data-item occurrence identity, and production versus consumption direction. Reviewed approaches can represent these relationships, but admissible records do not uniformly require all four answers. The review did not identify a shared contract that both requires them and provides a common basis for comparison.

TIDE represents the selected distinctions through Execution and State identity carriers with input and output incidence. Once the carriers are fixed, incidence determines the selected relationship answers; the erasure witnesses show what can be lost when each distinction is removed. Structural and semantic validity constrain representation form, production attribution, and execution ancestry. Under supplied occurrence maps, preservation and reflection compare mapped incidence while outside and absent information remain separate.

The Dagster and OpenLineage projections have complete agreement under the supplied correspondence; controlled cases separate scope, mapped disagreement, absence, and semantic invalidity. These results establish what the comparator determines after projection and correspondence are supplied; they do not establish that a supplied projection is true or complete or that supplied correspondence is correct.

\section*{Acknowledgment}
ChatGPT (OpenAI)~\cite{chatgpt} supported drafting and restructuring
throughout Sections I--IX, source checking, and artifact analysis.
The author originated and approved the model, questions, scope,
evidence, and conclusions, verified the manuscript, and accepts
responsibility for every claim; ChatGPT is not an author.

\begin{thebibliography}{15}
\setlength{\itemsep}{0pt}
\setlength{\parskip}{0pt}

\bibitem{openlineage-object}
OpenLineage, ``Object model,'' \emph{OpenLineage Core Specification}, accessed Aug. 6, 2026. [Online]. Available: \url{https://openlineage.io/docs/spec/object-model/}

\bibitem{mlmd-guide}
TensorFlow, ``ML Metadata,'' accessed Aug. 2, 2026. [Online]. Available: \url{https://www.tensorflow.org/tfx/guide/mlmd}

\bibitem{wfrun}
S. Leo \emph{et al.}, ``Recording provenance of workflow runs with RO-Crate,''
\emph{PLOS ONE}, vol. 19, no. 9, e0309210, 2024, doi: 10.1371/journal.pone.0309210.

\bibitem{mediator}
T. Ellkvist, D. Koop, J. Freire, C. T. Silva, and L. Str\"omb\"ack,
``Using mediation to achieve provenance interoperability,''
in \emph{Proc. IEEE Congress on Services-I}, 2009, pp. 291--298, doi: 10.1109/SERVICES-I.2009.68.

\bibitem{ppif}
A. Prabhune, A. Zweig, R. Stotzka, J. Hesser, and M. Gertz,
``P-PIF: A ProvONE provenance interoperability framework for analyzing heterogeneous workflow specifications and provenance traces,''
\emph{Distrib. Parallel Databases}, vol. 36, no. 1, pp. 219--264, 2018, doi: 10.1007/s10619-017-7216-y.

\bibitem{artifact}
    G. Dolbear, ``TIDE exact-relationship conformance and question-audit artifact,'' 2026. [Online]. Available: \url{https://github.com/tnkinobirb/TIDE-Model-Artifact}

\bibitem{second-challenge}
Provenance Challenge, ``Second Provenance Challenge,'' accessed Aug. 3, 2026. [Online]. Available: \url{https://openprovenance.org/provenance-challenge/SecondProvenanceChallenge.html}

\bibitem{third-challenge-questions}
Provenance Challenge, ``Provenance Challenge 3 questions,'' accessed Aug. 2, 2026. [Online]. Available: \url{https://openprovenance.org/provenance-challenge/ProvenanceQuestionsPc3.html}

\bibitem{mlmd-tutorial}
TensorFlow, ``Better ML engineering with ML Metadata,'' accessed Aug. 2, 2026. [Online]. Available: \url{https://www.tensorflow.org/tfx/tutorials/mlmd/mlmd_tutorial}

\bibitem{opm}
L. Moreau \emph{et al.}, ``The Open Provenance Model core specification (v1.1),''
\emph{Future Gener. Comput. Syst.}, vol. 27, no. 6, pp. 743--756, 2011, doi: 10.1016/j.future.2010.07.005.

\bibitem{prov-dm}
L. Moreau and P. Missier, Eds., ``PROV-DM: The PROV data model,''
W3C Recommendation, Apr. 2013. [Online]. Available: \url{https://www.w3.org/TR/prov-dm/}

\bibitem{prov-constraints}
J. Cheney, P. Missier, and L. Moreau, ``Constraints of the PROV data model,''
W3C Recommendation, Apr. 2013. [Online]. Available: \url{https://www.w3.org/TR/prov-constraints/}

\bibitem{cpf}
U. J. Braun, M. I. Seltzer, A. Chapman, B. Blaustein, M. D. Allen, and L. Seligman,
``Towards query interoperability: PASSing PLUS,''
in \emph{Proc. 2nd USENIX Workshop on the Theory and Practice of Provenance (TaPP)}, 2010. [Online]. Available: \url{https://www.usenix.org/conference/tapp-10/towards-query-interoperability-passing-plus}

\bibitem{wfrun-profile}
Workflow Run RO-Crate Working Group, ``Process Run Crate, version 0.5,'' accessed Aug. 5, 2026. [Online]. Available: \url{https://w3id.org/ro/wfrun/process/0.5}

\bibitem{provabs}
P. Missier, J. Bryans, C. Gamble, V. Curcin, and R. Danger,
``ProvAbs: Model, policy, and tooling for abstracting PROV graphs,''
in \emph{Provenance and Annotation of Data and Processes}, 2015, pp. 3--15, doi: 10.1007/978-3-319-16462-5\_1.

\bibitem{chatgpt}
OpenAI, ``ChatGPT,'' accessed Aug. 7, 2026. [Online].
Available: \url{https://chatgpt.com/}

\end{thebibliography}
\end{document}
